\documentclass[11pt,a4paper]{article}
\usepackage[hyperref]{acl2023}
\usepackage{times}
\usepackage{latexsym}
\usepackage{graphicx}
\usepackage{booktabs}
\usepackage{amsmath}

\title{Does Translation Preserve the Emotional Arc? \\
Comparing Human and LLM Literary Translation Fidelity}

\author{
  Julie Rhee \\
  Stanford University \\
  \texttt{juliarhee@stanford.edu}
}

\begin{document}
\maketitle

\noindent \textbf{Project type:} Custom \\
\textbf{TA Mentor:} [Your TA name here]

\begin{abstract}
Literary translation must preserve not just semantic content but the emotional trajectory of a narrative.
We investigate whether LLM-generated translations preserve emotional arcs more faithfully than established human translations, using five French--English novels from Project Gutenberg.
We construct a pipeline that splits texts into chapters, generates LLM translations via Gemini 2.5 Pro, and compares emotional arcs using VADER, labMT, and XLM-RoBERTa sentiment analysis.
Preliminary results on \textit{Wuthering Heights} (English$\to$French) show that the LLM translation achieves a significantly higher Pearson correlation with the original's emotional arc ($r{=}0.49$, $p{<}0.01$) than the human translation ($r{=}0.34$, $p{=}0.07$) under chapter-level XLM-RoBERTa analysis, suggesting LLMs may better preserve emotional fidelity in translation.
\end{abstract}

\section{Approach}

We frame translation fidelity as \textit{emotional arc preservation}: a faithful translation should produce a sentiment trajectory that closely tracks the original.
Our pipeline has four stages.

\paragraph{Corpus and Translation.}
We use five novels spanning both translation directions:
four French originals with English human translations (Voltaire's \textit{Candide}, Flaubert's \textit{Madame Bovary}, Verne's \textit{Twenty Thousand Leagues} and \textit{Around the World in 80 Days}), plus Bront\"{e}'s \textit{Wuthering Heights} (English original, French human translation).
All source texts and human translations are from Project Gutenberg.
For each book, we generate an LLM translation using Gemini 2.5 Pro on Google Cloud Vertex AI, translating chapter-by-chapter with a literary translation system prompt (temperature 0.3).
This yields three versions per book: the original, the human translation, and the LLM translation.

\paragraph{Sentiment Analysis.}
We score each chapter with three complementary methods:
\begin{itemize}
\setlength\itemsep{0pt}
    \item \textbf{VADER} \cite{hutto2014vader}: English-only lexicon+rule-based sentiment (compound score).
    \item \textbf{labMT} \cite{dodds2011temporal}: French-only happiness lexicon from hedonometer.org, rescaled from $[1,9]$ to $[-1,1]$.
    \item \textbf{XLM-RoBERTa} \cite{barbieri2022xlmtwitter}: Multilingual transformer fine-tuned on sentiment (\texttt{cardiffnlp/twitter-xlm-roberta-base-sentiment-multilingual}). We chunk chapters into ${\sim}400$-token windows with 50-token overlap, map predictions to $[-1,1]$, and take a weighted average.
\end{itemize}
XLM-RoBERTa is our primary cross-lingual method as it applies to both languages.
We also implement a sliding-window variant following \citet{reagan2016emotional}, sampling 100 evenly-spaced windows of 10{,}000 words across the full text to produce smooth, position-normalized arcs.

\paragraph{Arc Comparison Metrics.}
For each pair of versions (original vs.\ human, original vs.\ LLM, human vs.\ LLM), we compute:
Pearson and Spearman correlation, Dynamic Time Warping (DTW) distance normalized by path length, RMSE, mean absolute difference, cosine similarity, and peak/trough alignment (average positional shift of top-$k$ extrema as \% of arc length).
We compare arcs at both chapter-level and percentage-normalized (100 points) granularity.

\section{Experiments}

\paragraph{Data.}
Our corpus consists of five novels totaling ${\sim}150$ chapters.
All texts are sourced from Project Gutenberg.
\textit{Wuthering Heights} (34 English chapters, 30 French chapters in the human translation) is fully processed; the remaining four books require LLM translation in the next phase.

\paragraph{Evaluation.}
Our primary evaluation metric is Pearson correlation between sentiment arcs, which captures whether the emotional trajectory \textit{shape} is preserved.
DTW distance captures temporal misalignment, while peak/trough alignment measures whether key emotional moments occur at corresponding narrative positions.

\paragraph{Experimental Details.}
LLM translations were generated using Gemini 2.5 Pro via Google Cloud Vertex AI (\texttt{temperature=0.3}, \texttt{max\_output\_tokens=65536}).
Translation of \textit{Wuthering Heights} (34 chapters) consumed 165{,}961 input tokens and 199{,}127 output tokens.
XLM-RoBERTa inference uses the HuggingFace \texttt{transformers} library with 400-token chunks.
All code is original except for the pre-trained models cited above.

\paragraph{Results.}
Table~\ref{tab:results} reports chapter-level XLM-RoBERTa results for \textit{Wuthering Heights}.
The LLM translation shows a statistically significant correlation with the original ($r{=}0.49$, $p{=}0.008$), while the human translation's correlation is weaker and marginally significant ($r{=}0.34$, $p{=}0.069$).
The human--LLM pair shows low correlation ($r{=}0.19$), indicating the two translations diverge substantially from each other despite both tracking the original to varying degrees.

\begin{table}[h]
\centering
\small
\begin{tabular}{lccc}
\toprule
\textbf{Pair} & \textbf{Pearson $r$} & \textbf{$p$-value} & \textbf{Cosine} \\
\midrule
Orig vs Human  & 0.337 & 0.069 & 0.882 \\
Orig vs LLM    & \textbf{0.485} & \textbf{0.008} & 0.910 \\
Human vs LLM   & 0.192 & 0.317 & 0.944 \\
\bottomrule
\end{tabular}
\caption{Chapter-level XLM-RoBERTa arc comparison for \textit{Wuthering Heights} (EN$\to$FR). Bold indicates statistically significant results ($p < 0.05$).}
\label{tab:results}
\end{table}

The sliding-window analysis (100 points, 10{,}000-word windows) shows weaker correlations overall, with the human--LLM pair achieving the highest correlation ($r{=}0.34$, $p{<}0.001$), suggesting that at finer granularity the two French translations track each other more closely than either tracks the English original.
This may reflect systematic differences in how French and English encode sentiment at the lexical level.

These preliminary results are encouraging: the LLM translation appears to preserve the chapter-level emotional arc of \textit{Wuthering Heights} more faithfully than the 19th-century human translation.
However, we note this is a single book in one direction (EN$\to$FR); the pattern may differ for FR$\to$EN translations where the human translations are into English.

\section{Future Work}

\begin{itemize}
\setlength\itemsep{0pt}
    \item \textbf{Complete corpus:} Run LLM translations for the four remaining French$\to$English books to test whether the pattern holds across translation directions and literary styles.
    \item \textbf{LLM judge evaluation:} Use Gemini as a sentiment judge to provide a fourth scoring method not reliant on pre-trained lexicons or models.
    \item \textbf{Qualitative analysis:} Our pipeline includes a divergence-detection module that extracts the most emotionally divergent passages between versions. We will use this for close-reading analysis of \textit{where} and \textit{why} translations diverge.
    \item \textbf{Statistical robustness:} Apply bootstrap confidence intervals to correlation estimates and test sensitivity to chunking parameters.
    \item \textbf{Multiple LLMs:} Compare translations from different LLMs (e.g., Gemini vs.\ Claude vs.\ GPT-4) to assess whether emotional fidelity varies by model.
\end{itemize}

\bibliographystyle{acl_natbib}
\begin{thebibliography}{5}
\bibitem[\protect\citename{Hutto and Gilbert}2014]{hutto2014vader}
C.~J. Hutto and Eric Gilbert.
\newblock 2014.
\newblock VADER: A parsimonious rule-based model for sentiment analysis of social media text.
\newblock In \emph{Proceedings of ICWSM}.

\bibitem[\protect\citename{Dodds et~al.}2011]{dodds2011temporal}
Peter~Sheridan Dodds, Kameron~Decker Harris, Isabel~M. Kloumann, Catherine~A. Bliss, and Christopher~M. Danforth.
\newblock 2011.
\newblock Temporal patterns of happiness and information in a global social network: Hedonometrics and Twitter.
\newblock \emph{PLoS ONE}, 6(12):e26752.

\bibitem[\protect\citename{Barbieri et~al.}2022]{barbieri2022xlmtwitter}
Francesco Barbieri, Luis Espinosa~Anke, and Jose Camacho-Collados.
\newblock 2022.
\newblock XLM-T: Multilingual language models in Twitter for sentiment analysis and beyond.
\newblock In \emph{Proceedings of LREC}.

\bibitem[\protect\citename{Reagan et~al.}2016]{reagan2016emotional}
Andrew~J. Reagan, Lewis Mitchell, Dilan Kiley, Christopher~M. Danforth, and Peter~Sheridan Dodds.
\newblock 2016.
\newblock The emotional arcs of stories are dominated by six basic shapes.
\newblock \emph{EPJ Data Science}, 5(1):31.

\bibitem[\protect\citename{Bront\"{e}}1847]{bronte1847wuthering}
Emily Bront\"{e}.
\newblock 1847.
\newblock \emph{Wuthering Heights}.
\newblock Thomas Cautley Newby, London.
\end{thebibliography}

\end{document}
